\section{Fourier--Bohr scalarization of the residual label}
\label{sec:tpc46-bohr}

The multiplicative label \(s=d(m)r\) can be retained without paying a
dimension factor.  This shows that the bare residual product is not the
source of the tensor-normalization loss.

Let \(\cQ\) be a finite prime set containing all prime factors of every
\(r\) and \(d(m)\) in the packet, and put
\begin{equation}
 \T^{\cQ}:=\prod_{q\in\cQ}\{z_q\in\C:|z_q|=1\}
 \label{eq:tpc46-Bohr-torus}
\end{equation}
with normalized Haar measure \(dz\).  For
\(z=(z_q)_{q\in\cQ}\), define the completely multiplicative character
\begin{equation}
 \chi_z(n):=\prod_{q\in\cQ}z_q^{v_q(n)}.
 \label{eq:tpc46-Bohr-character}
\end{equation}
Orthogonality gives
\begin{equation}
 \int_{\T^{\cQ}}\chi_z(s)\overline{\chi_z(t)}\,dz
 =\one_{s=t}
 \label{eq:tpc46-Bohr-orthogonality}
\end{equation}
for the supported integers \(s,t\).

\begin{theorem}[Residual-label fixed-shift closure]
\label{thm:tpc46-residual-fixed-shift}
For the residual packet in \eqref{eq:tpc46-residual-packet},
\begin{equation}
 \boxed{
 \sum_s|F(h,s)|^2
 \le
 \rho_{\cP,\cR}\rho_{\cM,\cJ}\cD_\otimes.}
 \label{eq:tpc46-residual-fixed-shift-bound}
\end{equation}
\end{theorem}

\begin{proof}
The Bohr transform of \(F(h,\cdot)\) is
\begin{align}
 \widetilde F_h(z)
 &:=
 \sum_sF(h,s)\chi_z(s)\notag\\
 &=\sum_{pr-mj=h}
 a_pb_rc_m\eta_j\chi_z(r)\chi_z(d(m)).
 \label{eq:tpc46-residual-Bohr-transform}
\end{align}
For each fixed \(z\), this is the scalar determinant form with
unit-modulus twists absorbed into \(b_r\) and \(c_m\).  Theorem
\ref{thm:tpc46-scalar-tensor-closure} therefore bounds
\(|\widetilde F_h(z)|^2\) by the right side of
\eqref{eq:tpc46-residual-fixed-shift-bound}, uniformly in \(z\).
Parseval in \eqref{eq:tpc46-Bohr-orthogonality} and integration over
\(z\) finish the proof.
\end{proof}

There is also a two-spectrum identity which records every remaining
collision exactly.  Put
\begin{align}
 A_z(\alpha)
 &:=
 \sum_{p,r}a_pb_r\chi_z(r)\e(-pr\alpha),
 \label{eq:tpc46-Az}\\
 C_z(\alpha)
 &:=
 \sum_{m,j}c_m\eta_j\chi_z(d(m))\e(mj\alpha).
 \label{eq:tpc46-Cz}
\end{align}

\begin{theorem}[Additive--multiplicative Parseval normal form]
\label{thm:tpc46-double-Parseval}
One has the exact identity
\begin{equation}
 \boxed{
 \sum_{h,s}|F(h,s)|^2
 =\int_{\T}\int_{\T^{\cQ}}
 |A_z(\alpha)|^2|C_z(\alpha)|^2\,dz\,d\alpha.}
 \label{eq:tpc46-double-Parseval}
\end{equation}
\end{theorem}

\begin{proof}
The joint Fourier--Bohr transform of \(F(h,s)\) equals
\begin{align*}
 \sum_{h,s}F(h,s)\chi_z(s)\e(-h\alpha)
 &=\sum_{p,r,m,j}a_pb_rc_m\eta_j
   \chi_z(r)\chi_z(d(m))
   \e(-(pr-mj)\alpha)\\
 &=A_z(\alpha)C_z(\alpha).
\end{align*}
Parseval on \(\T\times\T^{\cQ}\) gives
\eqref{eq:tpc46-double-Parseval}.
\end{proof}

Expansion of \eqref{eq:tpc46-double-Parseval} shows that its
off-diagonal terms are exactly the pairs satisfying both
\begin{equation}
 pr-mj=p'r'-m'j',
 \qquad
 d(m)r=d(m')r'.
 \label{eq:tpc46-mixed-collisions}
\end{equation}
Thus the remaining Gram problem is a mixed additive--multiplicative
energy, not an unspecified appeal to M\"obius randomness.

For squarefree-kernel labels, the same argument uses the finite Walsh
group \(\{\pm1\}^{\cQ}\) and prime valuations modulo two.  At the
physical equality alias, TPC coprimality makes the simpler product
version above literal.
